% ============================================================
% Chapter 10 --- Learning the innovation law: EM for non-Gaussian chains
%
% Original work of this thesis. Builds directly on the grid-BP machinery
% of Chapter~\ref{ch:laplace} (\ref{sec:l-grid}), where the innovation law
% was assumed known. Here it is estimated. Numerical values are those of
% research/nongaussian-bp; full derivations and the complete experiment
% suite (27 experiments) are in that package's compendium and paper.
% ============================================================

\chapter{Learning the Innovation Law: EM for Non-Gaussian Chains}
\label{ch:em-method}

Chapters~\ref{ch:gaussian} and~\ref{ch:laplace} treated the transition law of
the clean chain as given: Gaussian in the first case, a fixed non-Gaussian
family in the second. Belief propagation then computed the diffusion score
exactly, in the Gaussian case, or against a validated numerical reference, in
the Laplace case. Neither chapter asked where the transition law itself comes
from. This chapter does. The clean data generating process is the same
first-order Markov chain of \eqref{eq:l-chain-prior}, but the innovation law
is now a member of an estimated family, and the question is how much the
Markov structure this thesis has exploited throughout still helps once part of
the model must be learned rather than assumed, and what breaks first when the
assumptions themselves are wrong.

\begin{quote}
How much does known Markov structure reduce the statistical and computational
burden of learning a diffusion score, and what is lost when the transition law
or the continuous messages must be approximated?
\end{quote}

The chapter answers the \emph{how} --- the estimator, what is exact in it and
what is numerical, and what the implementation was checked against.
Chapter~\ref{ch:em-results} answers the \emph{how much}.

\section{Model: an estimated mixture innovation}\label{sec:em-model}

The clean chain is the one of \eqref{eq:l-chain-prior}, in the index
convention of this chapter,
\begin{equation}
  a_1 \sim \Nd(0,1), \qquad a_i = \alpha\, a_{i-1} + \varepsilon_i,
  \qquad i = 2,\dots,n,
  \label{eq:em-chain}
\end{equation}
with innovations $\varepsilon_i$ i.i.d., mean zero, variance
$\ivar = 1-\alpha^2$ --- the same Gaussian initial state and the same
variance-matching that Chapter~\ref{ch:laplace} used --- but now drawn from
a member of a \emph{parametric family} rather than a fixed law. Forward noising is the coordinatewise
Ornstein--Uhlenbeck channel of Section~\ref{sec:ou},
$x = e^{-t} a + \sqrt{\Deltat}\,\xi$, unchanged from every earlier chapter.

The estimated family is a $C$-component Gaussian mixture on the innovation.
It is the transition kernel written $M(a' \mid a)$ in
Chapters~\ref{ch:background} and~\ref{ch:laplace}, carrying its parameters
explicitly because they are now what is being estimated:
\begin{equation}
  K_\theta(a' \mid a) = \sum_{c=1}^{C} \pi_c \,
    \Nd\!\big(a' - \alpha a;\ \nu_c,\ s_c^2\big),
  \qquad \theta = \{\alpha, (\pi_c,\nu_c,s_c)_{c=1}^C\},
  \label{eq:em-mixture-kernel}
\end{equation}
so both the autoregressive coefficient $\alpha$ and the shape of the
innovation density are free parameters, estimated jointly by maximum marginal
likelihood from noised data alone. At $C=1$ this collapses to a single
Gaussian, recovering the model of Chapter~\ref{ch:gaussian}; larger $C$ can
represent skew, heavy tails, and multimodality to the resolution of the
grid. The initial law $a_1 \sim \Nd(0,1)$ is \emph{held fixed} at
estimation time --- it is the value used to generate every family compared in
Chapter~\ref{ch:em-results}, so there is no misspecification in it, but
neither is there direct evidence it could itself be recovered from noised
data; that question is not addressed here.

Two assumptions carry the argument and, consistent with the practice of this
thesis, are stated rather than tested by the estimator itself: the clean
process is exactly first-order Markov, and the forward noise is exactly
coordinatewise. Section~\ref{sec:em-nonmarkov} in the next chapter measures
what happens when the first of these is false by a controlled amount; the
second is never relaxed anywhere in this thesis.

\begin{remark}[Covariance-stationary, not strictly stationary]
Choosing $\ivar=1-\alpha^2$ makes the chain covariance-stationary --- $\operatorname{Var}(a_i)=1$
and $\operatorname{Cov}(a_i,a_j)=\alpha^{|i-j|}$ hold exactly for \emph{every} innovation
law, by the same second-moment recursion as the Gaussian case. It does not
make the chain strictly stationary: for a non-Gaussian innovation the
invariant marginal law of the chain is not $\Nd(0,1)$, so fixing $a_1\sim
\Nd(0,1)$ starts every family off its own invariant distribution. Nothing
compared in this or the next chapter depends on the distinction, because every
family compared shares the same $a_1$ and the same covariance structure. The
claim available here is covariance-stationarity, and it is the one used.
\end{remark}

\section{What remains exact}\label{sec:em-exact}

Two facts from Section~\ref{sec:l-grid} carry over unchanged, and everything
the estimator does rests on them, so they are worth restating here.

\begin{proposition}[The posterior factor graph is still the prior's chain]
\label{prop:em-unary}
Coordinatewise forward noising contributes exactly one \emph{unary} factor per
site, $x_i \mapsto \phi_i(a_i) \propto \exp\!\big(-(x_i-e^{-t}a_i)^2/2\Deltat\big)$.
It reweights node potentials and creates no new edges among the latent
$a_i$, for any innovation law, mixture or otherwise. The posterior over the
clean chain given the noised observation therefore factorises on exactly the
same chain graph as the prior.
\end{proposition}

\begin{proposition}[The E-step is exact]\label{prop:em-estep}
Since the posterior graph is a tree, one forward and one backward
sum--product sweep return the exact single-site and pairwise marginals of the
posterior, at the level of functional messages, for \emph{any} innovation
law. Consequently the E-step of an EM procedure for $\theta$ requires no
variational approximation and no sampling: the sufficient statistics EM needs
are exact conditional expectations under the true posterior, mixture label
included.
\end{proposition}

Both are consequences of the chain being a tree (Theorem~\ref{thm:tree-exact})
and hold independently of whether the innovation law admits a closed-form
message update. What is \emph{not} exact is the representation of a functional
message on a computer, which is where Section~\ref{sec:em-numerical} and the
grid-BP machinery of Section~\ref{sec:l-grid} enter.

\section{The EM algorithm}\label{sec:em-algorithm}

Write $\Xi(\theta)$ for the vector of exact posterior sufficient statistics
that BP supplies given the current parameter $\theta$ --- the conditional
expectations of the pairwise sums, mixture-label indicators, and second
moments that the mixture kernel's complete-data log-likelihood needs. The
gradient of the \emph{marginal} log-likelihood $\mathcal L(\theta)$ is given
by Fisher's identity,
\begin{equation}
  \nabla_\theta \mathcal L(\theta)
  = \big\langle \Xi(\theta),\ \nabla_\theta \log K_\theta \big\rangle .
  \label{eq:em-fisher}
\end{equation}
The practical content of \eqref{eq:em-fisher} is that belief propagation is
never differentiated through: BP supplies the conditional expectation
$\Xi(\theta)$ at the current parameter, and the gradient with respect to
$\theta$ is then an ordinary derivative of the kernel's log-density. This
is what makes the E-step of Proposition~\ref{prop:em-estep} directly usable
inside an optimiser.

The M-step for the mixture kernel is \emph{not} an exact maximisation. Given
$\Xi(\theta^{(k)})$, updating $(\pi_c,\nu_c,s_c)_{c=1}^C$ exactly would
require solving a mixture-likelihood problem with its own latent label, so
the implementation alternates closed-form \emph{conditional} maximisers ---
the mixture block and $\alpha$ --- until the per-edge gain in $Q$ falls below
a tolerance or a cap of sixteen sweeps is reached, recording on the returned
kernel which of the two occurred. The algorithm is therefore generalised EM,
or ECM \citep{meng1993maximum}, rather than exact EM: a sweep ascends $Q$
without maximising it.

The cap of sixteen is set by measurement, not convention. Varying only this
budget at $C=8$, $n=256$, and a fixed outer count of $40$, the fitted
innovation excess kurtosis moves by $28\%$ between four sweeps and sixteen
($1.68 \to 2.14$, against a true value of $3.0$) while $\alpha$ moves in the
fourth decimal --- the same rate asymmetry between shape and correlation that
Section~\ref{sec:em-convergence} finds in the outer loop, appearing one level
down. The effect is confined to short outer runs: by $60$ outer iterations
the two loops substitute for one another and the inner budget stops
mattering. Sixteen removes the interaction rather than requiring a reader to
establish which regime a given cell is in.

Monotone ascent of $\mathcal L$ is guaranteed by the standard EM argument
applied to the surrogate $Q(\theta \mid \theta^{(k)}) = \langle
\Xi(\theta^{(k)}), \log K_\theta\rangle$ --- ascending $Q$ is all
monotonicity requires --- and is checked directly: across an executable audit
of the implementation, the minimum increase in $Q$ over twelve consecutive
steps was $+2.61$ for a Gaussian kernel, $+4.34$ for Laplace, and $+4.18$
for an eight-component mixture, with zero decreases recorded in any logged
run. The returned kernel's evidence equals
\texttt{trace.log\_evidence[-1]} by construction, with a regression test
asserting it, so the reported convergence trace always describes the model
actually returned.

\section{The numerical representation: grid BP, again}\label{sec:em-numerical}

Messages remain functions on $\R$, so they are represented exactly as in
Section~\ref{sec:l-grid}: values on $\ngrid$ equispaced points in $[-A,A]$, updates
by trapezoidal quadrature. Two error sources are tracked separately here as
well. Truncation error --- the probability mass the finite domain $[-A,A]$
misses --- is diagnosed by the maximum boundary mass over the sites of a chain.
That is a statistic of a sampled trajectory rather than a bound, so it is
recorded over $256$ chains per configuration: the worst chain falls from
$2.3\times10^{-3}$ at $A=4$ to $1.2\times10^{-6}$ at $A=8$, the $90$th
percentile from $2.7\times10^{-4}$ to $1.5\times10^{-8}$, and neither depends on
the grid resolution $\ngrid$. The gap between the two rows matters, because a figure
read off a single trajectory sits near the median and so understates the tail by
two orders of magnitude. Quadrature error, from the
finite spacing $h=2A/(\ngrid-1)$ --- the grid holds $\ngrid$ points including both
endpoints --- on an otherwise exact interior update, falls as
$O(h^2)$: a factor of four per halving of $h$, matching the trapezoidal
rate on a smooth integrand. At the working configuration used for every
experiment in this and the next chapter, $\ngrid=401$, $A=8$, both sources sit far
below any effect subsequently measured: the discretised recursion reproduces
the closed-form Gaussian posterior mean to $9.2\times10^{-15}$, and agrees
with an independently coded brute-force enumeration --- sharing no code with
the BP recursion --- to $1.6\times10^{-14}$ on posterior means and
$1.8\times10^{-15}$ on the log-evidence.

\paragraph{Which finite-grid model is being optimised.} ``Exact for the
discretised chain'' needs one more specification than it has been given so far,
because there are two conventions available and they lead to different
objectives. The transition array used here is the continuum density
$K_\theta(a' \mid a)$ evaluated at the grid points and combined with the
trapezoidal weights, with \emph{no} per-column renormalisation. BP is therefore
exact --- up to floating point --- for the chain-shaped \emph{factor model} that
array defines, and that model is not literally a normalised finite-state
transition kernel on the truncated domain. The distance between the two is
precisely the column-mass residual already reported as the quadrature
diagnostic. At the working configuration, the worst column with $|a| \le A/2$
departs from unit mass by $\gdinterior$; such a column sits far enough inside
the domain that its transition density has no tail outside the grid, so its
residual measures quadrature alone. The worst column anywhere departs by
$\gdedge$, at the domain edge where the density is truncated.

The alternative convention --- renormalising each parent-conditioned column
under the quadrature weights --- would make the array a genuine finite-state
kernel. Its normaliser would then depend on $\theta$, most strongly near the
boundary, and that dependence would have to be carried: the derivative of
$\log Z_\theta$ would enter both Fisher's identity and the M-step, and the
surrogate could no longer be written as
$Q(\theta \mid \theta^{(k)}) = \langle \Xi(\theta^{(k)}), \log K_\theta\rangle$
with $K_\theta$ the raw density. Because no renormalisation is applied, that
term is absent by construction rather than by oversight, and the surrogate above
is the correct one for the model actually fitted. The price is that the fitted
object is the continuum kernel read through a quadrature rule, which is what the
column-mass residual measures and what the resolution gate of
Section~\ref{sec:em-numerical} bounds.

Figure~\ref{fig:em-grid-domain} separates the two error sources the paragraphs
above keep apart: resolution and truncation are refined one at a time, so that
neither diagnostic can be mistaken for a certificate of the other.

\begin{figure}[t]\centering
\includegraphics[width=1.0\linewidth]{fig_grid_domain.pdf}
\caption[Grid and domain control, separated]{Grid and domain control on a Laplace chain ($\alpha=0.85$, $\nsites=32$
sites), four panels. (a) Relative score error against a high-resolution
reference as grid resolution $\ngrid$ grows at fixed half-width $A=8$; (b) the same
as $A$ grows at fixed $\ngrid=401$ --- both saturate at the floating-point floor
well inside the working configuration. (c) Forward-message boundary mass
against $A$, the truncation diagnostic, drawn as \emph{both} the maximum and
the $90$th percentile over $256$ chains, which differ by roughly two orders of
magnitude for the reason given in Remark~\ref{rmk:em-boundary}. (d) Interior
column-mass residual against spacing at $A=8$, the quadrature diagnostic, with
an $O(h^2)$ reference. Neither diagnostic alone certifies convergence. Source:
\texttt{outputs/frozen/exp\_01/} and \texttt{outputs/frozen/exp\_18/}.}
\label{fig:em-grid-domain}
\end{figure}

\begin{remark}[The boundary-mass diagnostic is empirical, not a bound]
\label{rmk:em-boundary}
A boundary-mass figure computed from one sampled chain does not bound the
truncation error for every run, and the gap is large enough to matter. Running
the same recursion on one chain and on $\gdchains$ of them, at the working
configuration and over all $\gdlevels$ noise levels, the batch maximum exceeds
the single-chain reading by a median factor of $\gdboundmed$ and by as much as
$\gdboundhi$; at the mildest level it is only $\gdboundlo$. A single draw
therefore lands somewhere in the body of a distribution with a very long tail,
and where it lands is not knowable from the draw. This is why the figures above
are reported as a maximum and a $90$th percentile over chains rather than as one
number. The diagnostic is an empirical estimate of typical truncation
loss, not a certified bound. What it does \emph{not} affect is the separation of truncation from
quadrature error reported above, which is a pure function of $(A,\ngrid)$ and the
prior and involves no sampled chain at all --- that comparison reproduces
to $0.000$ relative difference under reseeding and is exactly the result
Section~\ref{sec:l-grid} already relies on.
\end{remark}

One further identity is worth restating from Section~\ref{sec:l-grid} because
it governs every result in the next chapter exactly as it governed
Chapter~\ref{ch:laplace}: whatever an approximation --- Gaussian closure, grid
truncation, or an under-converged fit --- does to the posterior mean is
amplified by $e^{-t}/\Deltat \sim 1/2t$ in the score, so small diffusion times
remain the regime where any deficiency is most visible.

Figure~\ref{fig:em-diagnostics} collects the optimisation-side checks:
monotone ascent of the marginal likelihood, recovery of $\alpha$ from
initialisations far from the truth, and the $n^{-1/2}$ decay of the recovery
error. The first two are what make the procedure an EM rather than a heuristic;
the third is what makes it an estimator.

\begin{figure}[t]\centering
\includegraphics[width=1.0\linewidth]{fig_em_diagnostics.pdf}
\caption[Optimisation diagnostics for the EM procedure]{Optimisation diagnostics for the EM procedure of
Section~\ref{sec:em-algorithm}. Left: marginal log-likelihood by iteration
from four initialisations of $\alpha$; the largest decrease across
iterations is zero in every run. Centre: the estimated $\alpha$ over the
same runs, converging to $[0.8517,0.8520]$ against a truth of $0.85$ from
initialisations spanning $[-0.4,0.6]$ --- direct evidence that $\alpha$ is
estimated, not supplied. Right: RMSE of the recovered $\alpha$ and
innovation variance against the number of training sequences, with an
$n^{-1/2}$ reference line, at the same truth of $0.85$ and over $16$ seeds.
Appendix~\ref{app:repro} maps every panel to the run that produced it.}
\label{fig:em-diagnostics}
\end{figure}

\section{What convergence actually costs}\label{sec:em-convergence}

The last piece of machinery this chapter reports is not a theorem but a
measurement, and it is the one that most directly qualifies the results of
Chapter~\ref{ch:em-results}: how many EM iterations the algorithm of
Section~\ref{sec:em-algorithm} actually needs.

Table~\ref{tab:em-convergence} shows what the channel costs, in two parts: the
through-channel fit at four updates, and where a clean and a channel arm each
finish.

%% GENERATED by tools/make_convergence_numbers.py -- do not hand-edit.
%% Trace: outputs/frozen/exp_27_seed16/shape_trace.csv (16 seeds). Endpoints: outputs/frozen/exp_06/clean_vs_noised_shape.csv.

\begin{center}
\begin{minipage}{\linewidth}\centering
\captionof{table}[What the channel costs in convergence rate]{What the
observation channel costs in rate rather than in where the fit lands. Above:
the through-channel fit at $\nseq = 512$, $C = 8$, as a
median over $16$ seeds at four updates ---
the correlation is essentially in place by $30$ while the innovation shape is
still moving at $120$. Below: where the two arms finish, fit on clean transition
pairs and on the same chains once through the channel. Generating values are
$\corr = 0.85$ and excess kurtosis $3.0$; a finite mixture does not contain the
generating law exactly, so what is approached is the best fit within the family.}
\label{tab:em-convergence}
\begin{tabular}{@{}lcc@{}}
\toprule
update & $\corr$ & innovation excess kurtosis \\
\midrule
$1$ & $0.4402$ & $0.078$ \\
$30$ & $0.8412$ & $1.275$ \\
$120$ & $0.8499$ & $3.019$ \\
$600$ & $0.8504$ & $2.980$ \\
\midrule
\multicolumn{3}{@{}l@{}}{\emph{where each arm finishes}} \\
clean transition pairs & $0.8509$ & $2.939$ \\
through the channel & $0.8526$ & $2.729$ \\
\bottomrule
\end{tabular}
\end{minipage}
\end{center}


\noindent The two halves say different things. The trace shows the coordinates
separating: the correlation is essentially in place by thirty updates while the
innovation shape is still climbing at a hundred and twenty, because the mixture
label is a latent variable of its own and the channel compounds the missing
information --- the mechanism of \citet{dempster1977maximum} appearing in the
convergence rate, where the theory places it.

The channel has two distinct effects, and they should not be collapsed into
one. First, it changes the local rate of EM through the fraction of missing
information, which is the effect the trace above displays. Second, at finite
sample size it reduces the observed Fisher information and can therefore
increase estimation uncertainty --- an effect measured directly, by coordinate,
in Section~\ref{sec:em-information}. Under identifiability and correct
specification it need not move the population optimum, and that is the limited
sense in which the clean and channel arms may approach the same endpoint while
doing so at different rates. The endpoints here are consistent with that
reading: the two arms finish within $0.002$ of each other on the correlation
and within a fifth of a unit on excess kurtosis, both close to the generating
values. This is evidence of similar fitted endpoints in this configuration, not
of an absent information or accuracy penalty.

That is one configuration. A dedicated sweep measured the cost across sixteen
seeds, recording for each fit the first update after which a coordinate stays
within a fixed relative tolerance of its value at the end of the run. At
$\nseq = \convsize$, $C = \convcomps$ and the \convdesign{} noise design, over
$\convseeds$ seeds run to $\convcap$ updates, the correlation coefficient
settles at a median of $\convrhomed$ updates and the innovation shape at
$\convshapemed$, the largest observed value being $\convshapemax$. Forming the
ratio within each fit and then taking the median over fits gives $\convratio$;
the quotient of the two medians is $\convratioofmedians$, and the paired figure
is the one reported, because both coordinates are measured on the same fit.

Two things about what that measures. ``Settled'' is trace stability --- the
first update after which a coordinate stays within $10^{-3}$ of its end-of-run
value for $\alpha$, or $2\times10^{-2}$ for the excess kurtosis --- and not
distance to an optimiser; a coordinate can sit still while the density it
belongs to moves along a shallow direction. And the run has to be long enough
for the statistic to mean anything: the sweep was run at two lengths, and these
medians come from the $\convcap$-update shard, the only one in which no
configuration settles near its own cap. Against them, the fixed budget of
$\convbudget$ iterations that several experiments in
Chapter~\ref{ch:em-results} used as a default
falls short of the shape coordinate's median by about a factor of six and
leaves $\convunsettled$ of the sweep's $\convunsettledof$ configurations ---
both shards, since censoring can only shorten a settling time and so cannot
inflate this count --- unsettled in it.

\begin{remark}[Fixed budgets compare convergence rates, not estimators]
\label{rmk:em-fixed-budget}
This is the single methodological lesson of the chapter, and it recurs
throughout the next one. Whatever coordinate of a fit converges slower looks
worse under a fixed iteration count, and a slower-converging configuration
under-iterated by the same budget produces what looks like a genuine deficit
rather than symmetric noise --- which is precisely the signature that makes it
easy to mistake for a real effect. Compared correctly, at eight seeds and a
fixed budget of 200 updates, no channel penalty is resolved in the reported
shape summary: the paired clean $-$ noised difference at that operating point
is $-0.052\pm0.217$, consistent with zero. The operating point is not a
converged one --- the chapter's own median settling time for the shape
coordinate is $\convshapemed$ updates --- so this is a statement about the
paired difference at $200$ updates, not about the converged estimators. A single non-converged draw would
not have supported that comparison in either direction. Where Chapter~\ref{ch:em-results}
reports a result obtained at a fixed, possibly insufficient budget, it says so
explicitly and states which quantities are and are not affected.
\end{remark}

\section{Summary}

The estimator combines two facts that hold for any innovation law ---
Proposition~\ref{prop:em-unary} and Proposition~\ref{prop:em-estep} --- with a
generalised M-step and Fisher's identity to turn exact BP conditional
expectations into a marginal-likelihood gradient. Everything about the
\emph{posterior inference} step is exact up to the grid discretisation of
Section~\ref{sec:em-numerical}, which is controlled to far below any effect
measured. What is not free is optimisation time. Observing the chain only
through the channel costs the autoregressive coefficient about an order of
magnitude in iterations, and the innovation shape then settles some three times
slower again, with a tail several times longer than its own median.
Chapter~\ref{ch:em-results} quantifies pointwise accuracy, capacity and
robustness to controlled misspecification, and leaves the reverse-generation
comparison open, because its fitted shape coordinates were not sufficiently
settled under the available budget.
